% Part: first-order-logic
% Chapter: models-theories
% Section: introduction

\documentclass[../../../include/open-logic-section]{subfiles}

\begin{document}

\olfileid{fol}{mat}{int}
\olsection{ভূমিকা}

\begin{explain}
স্বতঃসিদ্ধমূলক পদ্ধতির বিকাশ বিজ্ঞানের ইতিহাসে একটি উল্লেখযোগ্য অর্জন,
এবং গণিতের ইতিহাসে তার বিশেষ গুরুত্ব আছে। কোনো বিষয়ক্ষেত্রকে
স্বতঃসিদ্ধমূলকভাবে গড়ে তুলতে অনেকগুলি প্রশ্ন স্পষ্ট করতে হয়:
বিষয়ক্ষেত্রটি কী নিয়ে? তার সবচেয়ে মৌলিক ধারণাগুলি কী? সেগুলি পরস্পরের
সঙ্গে কীভাবে সম্পর্কিত? এই মৌলিক ধারণাগুলির সাহায্যে কি বিষয়ক্ষেত্রের
সব ধারণা সংজ্ঞায়িত করা যায়? এই ধারণাগুলি কোন বিধি মেনে চলে, এবং
অবশ্যই কোন বিধি মেনে চলতে হবে?

স্বতঃসিদ্ধমূলক পদ্ধতি ও যুক্তিবিদ্যা যেন পরস্পরের জন্যই তৈরি।
আনুষ্ঠানিক যুক্তিবিদ্যা স্বতঃসিদ্ধমূলক তত্ত্ব সূত্রবদ্ধ করার, তত্ত্বের
স্বতঃসিদ্ধগুলি থেকে সুনির্দিষ্ট পদ্ধতিতে উপপাদ্য প্রমাণ করার, এবং
স্বতঃসিদ্ধগুলি পরিতৃপ্ত করে এমন সব ব্যবস্থার ধর্ম নিয়মমাফিক অধ্যয়ন
করার উপকরণ দেয়।
\end{explain}

\begin{defn}
!!{sentence}s-এর কোনো সেট~$\Gamma$ \emph{বদ্ধ}, যদি এবং কেবল যদি
$\Gamma \Entails !A$ হলেই $!A \in \Gamma$ হয়। !!{sentence}s-এর
কোনো সেট~$\Gamma$-র \emph{আবরণ} হলো
$\Setabs{!A}{\Gamma \Entails !A}$।

$\Gamma$ যদি~$\Delta$-র আবরণ হয়, তবে বলি বাক্যের
সেট~$\Delta$ \emph{$\Gamma$-কে স্বতঃসিদ্ধায়িত করে}।
\end{defn}

\begin{explain}
কোনো স্বতঃসিদ্ধমূলক তত্ত্বকে তার স্বতঃসিদ্ধসমষ্টি~$\Delta$ যে
বাক্যসমষ্টিকে স্বতঃসিদ্ধায়িত করে, সেই সেট হিসেবে ভাবতে পারি।
অন্যভাবে বললে, যে বিজ্ঞানকে স্বতঃসিদ্ধমূলকভাবে গড়ে তুলে অধ্যয়ন করতে
চাই, তার মৌলিক ধারণাগুলির জন্য অ-যৌক্তিক সংকেতসমৃদ্ধ একটি প্রথম-ক্রমের
ভাষা এবং বিজ্ঞানের মৌলিক বিধিগুলি প্রকাশকারী !!{sentence}s-এর একটি সেট
থাকলে, স্বতঃসিদ্ধগুলি থেকে অনুসৃত ওই ভাষার সব !!{sentence}s দিয়ে
তত্ত্বটির প্রতিনিধিত্ব করা যায়। এতে যেমন মাত্র একটি মৌলিক ধারণা ও
সরল স্বতঃসিদ্ধ নিয়ে আংশিক ক্রমের তত্ত্বের মতো সহজ উদাহরণ পড়ে, তেমনি
নিউটনীয় বলবিদ্যার মতো জটিল তত্ত্বও পড়ে।

যেসব গুরুত্বপূর্ণ যৌক্তিক তথ্যের জন্য স্বতঃসিদ্ধমূলক পদ্ধতির এই
আনুষ্ঠানিক রূপ এত তাৎপর্যপূর্ণ, সেগুলি নিচে দেওয়া হলো। ধরি,
$\Gamma$ কোনো তত্ত্বের স্বতঃসিদ্ধ-পদ্ধতি, অর্থাৎ বাক্যের
একটি সেট।
\begin{enumerate}
\item কোনো স্বতঃসিদ্ধ-পদ্ধতি কখন !!{structure}s-এর অভিপ্রেত একটি
  শ্রেণিকে ঠিকভাবে ধারণ করে, তা নির্ভুলভাবে বলা যায়। অর্থাৎ,
  !!{structure}s-এর কোনো নির্দিষ্ট শ্রেণিতে আমাদের আগ্রহ থাকলে,
  স্বতঃসিদ্ধ-পদ্ধতি~$\Gamma$ সেই শ্রেণিকে সফলভাবে ধারণ করবে যদি এবং
  কেবল যদি !!{structure}s-গুলি ঠিক সেই~$\Struct M$-গুলি হয় যাদের জন্য
  $\Sat{M}{\Gamma}$।
\item এ ব্যাপারে আমরা ব্যর্থ হতে পারি, কারণ এমন~$\Struct M$ থাকতে পারে
  যার জন্য $\Sat{M}{\Gamma}$, কিন্তু $\Struct M$ আমাদের অভিপ্রেত
  !!{structure}s-এর একটি নয়। এতে এমন স্বতঃসিদ্ধ যোগ করার প্রয়োজন হতে
  পারে যা~$\Struct M$-এ সত্য নয়।
\item অন্তত এই অর্থে সফল হলে যে, সব অভিপ্রেত !!{structure}s-এ
  $\Gamma$ সত্য, তাহলে $\Gamma \Entails !A$ হলেই বাক্য~$!A$ সব অভিপ্রেত
  !!{structure}s-এ সত্য। তাই কোনো বাক্যকে স্বতঃসিদ্ধগুলি থেকে অনুসৃত
  দেখিয়েই সেটি সব অভিপ্রেত !!{structure}s-এ সত্য দেখাতে আমরা যৌক্তিক
  উপকরণ (যেমন !!{derivation} পদ্ধতি) ব্যবহার করতে পারি।
\item কখনও কোনো অভিপ্রেত !!{structure}s মনে রেখে শুরু করি না; বরং
  স্বতঃসিদ্ধগুলি থেকেই শুরু করি। কয়েকটি মৌলিক ধারণা নিয়ে চাই সেগুলি
  নির্দিষ্ট কিছু বিধি পরিতৃপ্ত করুক, এবং সেই বিধিগুলি একটি
  স্বতঃসিদ্ধ-পদ্ধতিতে সংকেতায়িত করি। প্রথমেই যাচাই করতে চাই,
  স্বতঃসিদ্ধগুলি যেন পরস্পরবিরোধী না হয়। পরস্পরবিরোধী হলে এই বিধি
  মানে এমন কোনো ধারণা থাকতে পারে না, আর আমরা একটি অসংলগ্ন তত্ত্ব
  দাঁড় করানোর চেষ্টা করেছি। $\Gamma$-র একটি মডেল খুঁজে দেখিয়ে এই
  বিপর্যয় ঘটে না বলে যাচাই করতে পারি। আমাদের তত্ত্বের মডেল থাকলে
  যৌক্তিক পদ্ধতিতে সেগুলি অনুসন্ধান করতে এবং মডেল নির্মাণ করতেও পারি।
\item স্বতঃসিদ্ধগুলির স্বাধীনতাও একটি গুরুত্বপূর্ণ প্রশ্ন। কোনো
  স্বতঃসিদ্ধ অন্যগুলির ফলশ্রুতি হয়ে অপ্রয়োজনীয় হতে পারে।
  $\Gamma$-র স্বতঃসিদ্ধ~$!A$ যে অপ্রয়োজনীয়, তা
  $\Gamma \setminus \{!A\} \Entails !A$ প্রমাণ করে দেখাতে পারি। আবার
  $(\Gamma \setminus \{!A\}) \cup \{\lnot !A\}$ পরিতৃপ্তিযোগ্য দেখিয়ে
  প্রমাণ করতে পারি যে, স্বতঃসিদ্ধটি অপ্রয়োজনীয় নয়। যেমন, এইভাবেই
  দেখানো হয়েছিল যে জ্যামিতির অন্য স্বতঃসিদ্ধগুলি থেকে সমান্তরাল
  স্বীকার্যটি স্বাধীন।
\item আরেকটি গুরুত্বপূর্ণ প্রশ্ন হলো কোনো তত্ত্বে ধারণাগুলির
  সংজ্ঞেয়তা। ভাষা নির্বাচন করলেই তত্ত্বের মডেলগুলি কী দিয়ে গঠিত হবে,
  তা নির্ধারিত হয়। কিন্তু কোনো তত্ত্বের প্রত্যেকটি দিক তার মডেলগুলিতে
  আলাদাভাবে উপস্থাপন করতে হয় না। যেমন, প্রত্যেক ক্রম~$\le$ একটি
  সংশ্লিষ্ট কঠোর ক্রম~$<$ নির্ধারণ করে—একটি দেওয়া থাকলে অন্যটি
  সংজ্ঞায়িত করা যায়। তাই এমন ক্রমযুক্ত কোনো তত্ত্বের মডেলে সংশ্লিষ্ট
  কঠোর ক্রমটিও \emph{আলাদাভাবে} থাকা আবশ্যক নয়। সাধারণভাবে কখন একটি
  সম্পর্ককে অন্য সম্পর্কগুলির সাহায্যে সংজ্ঞায়িত করা যায়? কখন অন্য
  সম্পর্কগুলির সাহায্যে কোনো সম্পর্ক সংজ্ঞায়িত করা অসম্ভব (এবং তাই
  তাকে ভাষার মৌলিক সংকেতগুলির মধ্যে যোগ করতেই হয়)?
\end{enumerate}
\end{explain}


\end{document}
